\scanpage{116}
\mspage{\unsure{142|146}}
\deleted{Soit $X$ un espace topologique, \deleted{dont chaque \unsure{pt} a un syst. fondamental de voisinages} \inserted{ayant une base d'ouverts} $U$ ayant les propriétés suivantes:

(i) $U$ rétrocompact dans $X$, et $U$ quasi-cpt

(ii) Toute \deleted{suite} famille $(E_i)_{i\in I}$ de parties \deleted{fermées} constructibles de $U$, recouvrant $U$, il y a une sous-famille finie recouvrant $U$.

Pour exprimer (ii), Introduisons [sur \unsure{tt} espace topologique $U$] la topologie dont les ouverts sont les \inserted{réunions quelconques de} parties loc. constructibles [c'est bien stable par intersections finies et réunions qlc] donc une base d'ouverts formée des parties constructibles. Lorsque $U$ est quasi-compact, et que \deleted{\unsure{tt pt admet}} admet une base d'ouverts rétrocompacts, ``loc. constructible'' est d'ailleurs pareil que ``constructible''; moyennant (i), les $U$ \unsure{retrouvent les} \unsure{(Hyp)} \unsure{donnant} dans le cas précédent.

NB. Si de plus $X$ satisfait l'axiome $T_0$, \inserted{[c'est le cas des préschémas]} \deleted{Alors} alors une telle topol. est \uline{séparée}, et même totalement discontinue [pour deux pts $x\ne y$, il existe des ouverts $V\ni x$, $W\ni y$, $V\cap W=\varnothing$, $V\cup W=U$]. L'Hyp (ii) exprime qu'elle est \uline{quasi-compacte} [donc \uline{compacte} dans le cas particulier que nous envisageons].}

\margin{\figure{\begin{tikzpicture}[x=1cm,y=1cm]
\draw(0,0)--(.04,-2.15);
\draw(0,-5.3)ellipse(.15 and .2);\fill(0,-5.3)circle(.025);\draw(-.18,-5.56)--(.22,-5.06);
\end{tikzpicture}}}
\margin{115}
